
\documentclass[12pt,a4paper]{article}

\usepackage[T2A]{fontenc}
\usepackage[utf8]{inputenc}
\usepackage[serbianc]{babel}

\usepackage{amsmath,amssymb,amsthm}
\usepackage{graphicx}
\usepackage{xcolor}
\usepackage{array}
\usepackage{enumitem}
\usepackage{wrapfig}
\usepackage{geometry}

\geometry{
    left=2.5cm,
    right=2.5cm,
    top=2.5cm,
    bottom=2.5cm
}


\newtheorem{definition}{Дефиниција}[section]
\newtheorem{theorem}{Теорема}[section]
\newtheorem{lemma}{Лема}[section]

\begin{document}


\begin{titlepage}
    \center

    {\Large Математички факултет\par}
    \vspace{0.5cm}
    {\Large Семинарски рад\par}
    \vspace{2cm}

    {\Huge\bfseries Сајбер безбедност\par}
    \vspace{2cm}

    {\Large Основни принципи заштите информационих система\par}

    \vfill

    \begin{flushright}
        \textbf{Студент:} Матиа Лазовић\\
        \textbf{Број Индекса:} 64/2025\\
        \textbf{Датум:} \today
    \end{flushright}

\end{titlepage}


\tableofcontents
\newpage

\section{Uvod}

Сајбер безбедност је део рачунарства који се бави тиме како да заштитимо компјутере, мреже и податке од неовлашћеног приступа, злоупотребе и разних напада. Како интернет и дигиталне услуге постају све доступнији, сајбер безбедност добија све већи значај.

Циљ сајбер безбедности је да се осигура \textbf{поверљивост},
\textit{интегритет} и доступност података. То су три главна принципа за безбедност информационих система.

\begin{wrapfigure}{l}{0.35\textwidth}
 \includegraphics[width=0.32\textwidth]{slika.png}
    \caption{Симболичан приказ сајбер безбедности}
   \label{fig:slika}
\end{wrapfigure}

Сајбер безбедност укључује многе технике, као што су шифровање, контрола приступа, заштитни зидови, антивируси и системи за откривање напада. Слика~\ref{fig:slika} приказује симболичан приказ заштите рачунара.

\section{Основни принципи сајбер безбедности}
\subsection{Дефиниција}
\begin{definition}
Сајбер безбедност је скуп метода, технологија и поступака којима се штите рачунарски системи, мреже и подаци од неовлашћеног приступа, уништавања, измене или крађе
\end{definition}

Заштита информационих система заснива се на неколико основних принципа.
Најважнији су:

\begin{itemize}
    \item поверљивост података;
    \item интегритет података;
    \item доступност система;
    \item аутентификација корисника;
    \item контрола приступа.
\end{itemize}

\subsection{Шифровање података}

Један од најважнијих начина заштите података је шифровање. Када се подаци шифрују, оригинална порука се помоћу кључа претвара у облик који неовлашћени корисници не могу да разумеју.

Једноставан математички модел може се представити формулом

\[
C=E_K(M),
\]

где је $М$ оригинална порука, $К$ кључ, а $C$ шифрована порука.

За дешифровање се користи одговарајући поступак:

\[
M=D_K(C).
\]

На тај начин корисник који поседује прави кључ може да добије оригиналну поруку.

\section{Врсте сајбер напада}

Постоји много различитих напада којима рачунарски системи могу бити угрожени.

\subsection{Фишинг}

Фишинг је напад где нападач покушава да превари корисника како би открио поверљиве информације, као што су лозинке или подаци о картицама. Напади се често изводе путем лажних имејлова и веб страница.

\subsection{Малвер}

Малвер је злонамерни софтвер направљен да поремети рад система, украде податке или омогући неовлашћени приступ рачунару.

\subsection{Напади на сервис}

Код ДДоС напада, нападач покушава да преоптерети одређени сервис великим бројем захтева. То може довести до привремене недоступности сервиса за легитимне кориснике.

\section{Математички аспект безбедности}

Математика игра веома важну улогу у модерној криптографији. Посебно су значајне теорија бројева, модуларна аритметика и вероватноћа.

На пример, модуларна аритметика може бити представљена изразом

\[
a \equiv b \pmod n.
\]

Овај концепт користи се у различитим криптографским алгоритмима.

\begin{lemma}
Ако су два цела броја једнака по модулу $n$, њихова разлика је дељива са $n$.
\end{lemma}

\begin{theorem}
Ако су $a$ и $b$ цели бројеви и важи $a \equiv b \pmod n$, тада постоји цео број $k$ такав да је $a-b=kn$.
\end{theorem}

\begin{proof}
По дефиницији конгруенције модуло $n$, израз $a \equiv b \pmod n$ значи да је $n$ делилац броја $a-b$. Зато постоји цео број $k$ за који важи

\[
a-b=kn.
\]

Тиме је тврђење доказано.
\end{proof}

\section{Поређење метода заштите}

У табели~\ref{tab:metode} риказане су неке од основних метода заштите рачунарских система.

\begin{table}[ht]
\centering

\label{tab:metode}
\begin{tabular}{|c|c|c|}
\hline
\textbf{Метода} & \textbf{Намена} & \textbf{Пример} \\
\hline
Шифровање & заштита података & AES \\
\hline
Аутентификација & провера идентитета & лозинка  \\
\hline
Заштитни зид & контрола мрежног саобраћаја & Firewall \\
\hline
Антивирус & откривање малвера & антивирусни програм \\
\hline
\end{tabular}
\caption{Методе заштите информационих система}
\end{table}

\section{Мере заштите}

Корисници могу значајно смањити ризик од сајбер напада применом основних безбедносних правила.

Нека основна правила:

\begin{itemize}
    \item користити јаке и јединствене лозинке;
    \item укључити вишефакторску аутентификацију;
    \item редовно ажурирати софтвер;
    \item правити резервне копије важних података;
    \item бити опрезан приликом отварања непознатих порука.
\end{itemize}
\begin{enumerate}
    \item идентификовати потенцијалне ризике;
    \item проценити ниво опасности;
    \item применити одговарајуће мере заштите;
    \item редовно проверавати безбедност система.
\end{enumerate}
    


Посебно је важно да корисници буду \textcolor{blue}{информисани} о могућим опасностима и да \textcolor{red}{не игноришу безбедносна упозорења}.

\section{Закључак}

Сајбер безбедност је један од најважнијих аспеката савременог дигиталног друштва. Заштита података и информационих система захтева комбинацију техничких решења, математичких метода и одговорног понашања корисника.
Примена шифровања, безбедне аутентификације, заштитних система и редовног ажурирања софтвера може значајно смањити ризик од сајбер напада. Због тога је развијање свести о сајбер безбедности важно и за појединце и за организације.

\end{document}